% Copyright 2026  Ed Bueler

\documentclass[10pt,
               svgnames,
               hyperref={colorlinks,citecolor=DeepPink4,linkcolor=FireBrick,urlcolor=Maroon},
               usepdftitle=false]{beamer}

\usetheme{Madrid}
\usecolortheme{beaver}
\setbeamercovered{transparent}
\setbeamerfont{frametitle}{size=\large}
\setbeamercolor*{block title}{bg=red!10}
\setbeamercolor*{block body}{bg=red!5}

\usepackage[english]{babel}
\usepackage[latin1]{inputenc}
\usepackage{times}
\usepackage[T1]{fontenc}
% Or whatever. Note that the encoding and the font should match. If T1
% does not look nice, try deleting the line with the fontenc.

\usepackage{empheq,bm,xspace}
%\usepackage{minted}
%\usepackage{hyperref}
\usepackage{tikz}

% If you wish to uncover everything in a step-wise fashion, uncomment
% the following command: 
%\beamerdefaultoverlayspecification{<+->}

\newcommand{\bb}{\mathbf{b}}
\newcommand{\bc}{\mathbf{c}}
\newcommand{\bbf}{\mathbf{f}}
\newcommand{\bl}{\bm{\ell}}
\newcommand{\br}{\mathbf{r}}
\newcommand{\bs}{\mathbf{s}}
\newcommand{\bx}{\mathbf{x}}
\newcommand{\by}{\mathbf{y}}
\newcommand{\bv}{\mathbf{v}}
\newcommand{\bu}{\mathbf{u}}
\newcommand{\bw}{\mathbf{w}}

\newcommand{\bzero}{\bm{0}}

\newcommand{\CC}{\mathbb{C}}
\newcommand{\RR}{\mathbb{R}}

\newcommand{\ddt}[1]{\ensuremath{\frac{\partial #1}{\partial t}}}
\newcommand{\ddx}[1]{\ensuremath{\frac{\partial #1}{\partial x}}}
\renewcommand{\t}[1]{\texttt{#1}}
\newcommand{\Matlab}{\textsc{Matlab}\xspace}
\newcommand{\eps}{\epsilon}

\newcommand{\twovect}[4]{\ensuremath{{#1}_{#2} =
                            \begin{bmatrix} #3 \\ #4 \end{bmatrix}}}

\newcommand{\rbullet}{{\color{FireBrick} \bullet}}

\newcommand*\circled[1]{\tikz[baseline=(char.base)]{
            \node[shape=circle,draw,inner sep=2pt] (char) {#1};}}


\title{Integration is a linear functional}

\subtitle{week 1 calculation}

\author{Ed Bueler}

\institute[]{UAF Math 617 Functional Analysis}

\date{Spring 2026}


\begin{document}
\beamertemplatenavigationsymbolsempty

\begin{frame}
  \maketitle
\end{frame}

\begin{frame}{Outline}
  \tableofcontents[hideallsubsections]
\end{frame}

\AtBeginSection[]
{% do not put toc here this time
}

\section{continuity and integration}


\begin{frame}{continuity and $C([a,b])$}

\begin{definition}[continuity]
\begin{itemize}
\item a function $f:[a,b] \to \RR$ is \emph{continuous at} $x\in[a,b]$ if for all $\eps>0$ there is $\delta>0$ so that
  $$|y-x|<\delta \text{ and } y\in[a,b] \quad \implies \quad |f(y)-f(x)|<\eps$$
\item a function $f:[a,b] \to \RR$ is \emph{continuous} if it is continuous at each $x\in[a,b]$
\end{itemize}
\end{definition}

\begin{definition}
   $$C([a,b]) = \left\{f:[a,b] \to \RR \,:\, f \text{ is continuous}\right\}$$
\end{definition}

\begin{itemize}
\item for example $a=0<1=b$:\quad  $C([0,1])$ \hfill [picture]

\medskip
\item note $C([0,1])$ is an infinite set

because it contains distinct

elements $1$, $x^1$, $x^2$, \dots
\vspace{20mm}
\end{itemize}

\end{frame}


\begin{frame}{$C([0,1])$ is vector space}

\begin{theorem}
$V=C([0,1])$ is a (real) vector space
\end{theorem}

\begin{proof}
If $f,g \in V$ then $h(x) = f(x)+g(x)=g(x)+f(x)$ is continuous,  so $V$ is closed under addition.  If $\alpha\in \RR$ then $\alpha f(x)$ is continuous, so $V$ is closed under scalar multiplication.   The zero function is a continuous function, and $f(x)+0=f(x)$, so $V$ has an additive identity.  Also $1 f(x) = f(x)$, so the scalar identity works.  And distribution rules are true: $\alpha(f+g) = \alpha f + \alpha g$, $\alpha(\beta f) = (\alpha \beta) f$.
\end{proof}

\vspace{50mm}
\end{frame}


\begin{frame}{definite integral}

\begin{definition}[Riemann's integral]
for a mesh $a = x_0 < x_1 < \dots < x_j < \dots < x_n = b$, and evaluation points $x_j^* \in [x_{j-1},x_j]$, define
  $$\int_a^b f(x)\,dx = \lim_{n\to\infty} \sum_{j=1}^n f(x_j^*)\,(x_j - x_{j-1})$$ 
\end{definition}

\begin{theorem}
if $f:[a,b] \to \RR$ is continuous then $\int_a^b f(x)\,dx \in \RR$ is well-defined
\end{theorem}

\begin{itemize}
\item this is \emph{not} the only definition of an integral
\item the limit process for the Riemann integral is more subtle than the limit process for continuity or derivatives
\item we will get to Lebesgue's integral later \dots and cover it lightly
\end{itemize}
\end{frame}


\begin{frame}{integration example}

Example: compute $\int_0^3 x\, dx$ from the definition

\vspace{60mm}
\end{frame}


\begin{frame}{fundamental theorem of calculus}

\begin{theorem}
if $f:[a,b] \to \RR$ is continuous and if $F:[a,b]\to \RR$ satisfies $F'(x)=f(x)$ then
    $$\int_a^b f(x)\, dx = F(b) - F(a)$$
\end{theorem}

\begin{itemize}
\item example: redo integral on last slide

\vspace{30mm}
\end{itemize}

\end{frame}


\begin{frame}{more integration examples}

example: $\displaystyle \int_0^1 x^k \,\frac{1}{1+x^2}\,dx$ for $k=0,1,2$

\vspace{60mm}
\end{frame}


\begin{frame}{more integration examples}

claim: for $k=0,1,2,\dots$,
$$\int_0^1 x^k \frac{1}{1+x^2}\,dx = c_k$$
where
\begin{align*}
c_k &= \frac{\pi}{4}, \frac{\ln(2)}{2}, 1 - \frac{\pi}{4}, \dots \\
   &= 0.78539, 0.34657, 0.21460, 0.15342, 0.11873, 0.09657, \dots
\end{align*}


\vspace{60mm}
\end{frame}


\AtBeginSection[] {
    \begin{frame}{Outline}
    \setbeamertemplate{section in toc}[sections numbered]
    \tableofcontents[currentsection]%,hideallsubsections]
    \end{frame}
}



\section{integration as a linear functional}

\begin{frame}{title}

\begin{itemize}
\item foo
   \begin{itemize}
   \item[$\circ$] bar
   \end{itemize}
\end{itemize}
\end{frame}


\section{numerical integration}

\begin{frame}{title}

\begin{itemize}
\item foo
   \begin{itemize}
   \item[$\circ$] bar
   \end{itemize}
\end{itemize}
\end{frame}


\section{normed vector spaces of functions}

\begin{frame}{title}

\begin{itemize}
\item foo
   \begin{itemize}
   \item[$\circ$] bar
   \end{itemize}
\end{itemize}
\end{frame}


\section{theory this week}

\begin{frame}{title}

\begin{itemize}
\item foo
   \begin{itemize}
   \item[$\circ$] bar
   \end{itemize}
\end{itemize}
\end{frame}



%\begin{frame}{references}
%
%\begin{itemize}
%{\small
%\item D.~P.~Kingma \& J.~Ba (2014). \href{https://arxiv.org/abs/1412.6980}{\emph{Adam: A method for stochastic optimization}}, preprint arXiv:1412.6980
%}
%\end{itemize}
%\end{frame}

\end{document}
